
This paper addresses the problem of reward signal design for GRPO-based code generation training. The central theoretical contribution is an analytical derivation showing that ratio rewards have 5--$20\times$ higher expected within-group variance than binary rewards in the partial-tractability regime (q $\in$ [0.3, 0.55], T = 5), with the variance ratio growing with the number of test cases T. This analysis provides a formal basis for preferring ratio rewards over binary rewards in GRPO settings where problems are partially solvable.

The prescreening protocol formalizes this insight into a computable criterion: the group tractability score S\_term and the target window S\_term $\in$ [0.3, 0.55]. A six-module prescreening pipeline implementing this protocol was validated at 15/15 tasks and 67/67 integration tests and is available for use by researchers applying GRPO to execution-based reward settings.

The primary experimental finding is negative: Qwen2.5-Coder-7B-Instruct achieves 0\% pass rate on 300 APPS introductory problems under a strict execution harness (2,400 total attempts), establishing that an SFT-initialized checkpoint is a prerequisite for any prescreened problems to exist. All mechanism experiments (h-m1 through h-m4) remain blocked. The empirical comparison of ratio versus binary reward under GRPO training --- whether ratio rewards produce more distinct advantage levels, higher gradient SNR, and earlier ZRF escape --- cannot be reported and is deferred to future work contingent on SFT checkpoint availability.

Three immediate next steps would complete the experimental program: (1) generate an SFT checkpoint by fine-tuning Qwen2.5-Coder-7B-Instruct on APPS introductory solutions; (2) run a format diagnostic to separate format incompatibility from capability failure as explanations for the 0\% result; (3) run a tractability validation with a stronger model to confirm that S\_term $\in$ [0.3, 0.55] is empirically achievable before committing SFT compute. Subject to these steps, the Stage 2 GRPO comparison requires no infrastructure changes and can proceed using the validated prescreening pipeline.

